% !TEX program = pdflatex
% Guide de prise en main - Portail RH Documents
% Compile: pdflatex ONBOARDING.tex
\documentclass[10.5pt,a4paper]{article}

\usepackage[T1]{fontenc}
\usepackage[utf8]{inputenc}
\usepackage[french]{babel}
\usepackage{microtype}
\usepackage{mathpazo}
\usepackage[scaled=0.92]{helvet}
\usepackage{courier}
\renewcommand{\familydefault}{\rmdefault}

\usepackage{geometry}
\usepackage{xcolor}
\usepackage{booktabs}
\usepackage{array}
\usepackage{tabularx}
\usepackage{colortbl}
\usepackage{ragged2e}
\usepackage{enumitem}
\usepackage{listings}
\usepackage{tikz}
\usepackage[explicit]{titlesec}
\usepackage{fancyhdr}
\usepackage[hidelinks]{hyperref}
\hypersetup{
  pdftitle={Guide de prise en main : Portail RH Documents},
  pdfauthor={Abdellah Raissouni},
}

\geometry{top=1.8cm, bottom=1.8cm, left=2.2cm, right=2.2cm,
          headheight=13pt, headsep=14pt}

\hyphenpenalty=10000
\exhyphenpenalty=10000
\tolerance=2000
\emergencystretch=2em
\AtBeginDocument{\sloppy}

% Brand palette
\definecolor{BrandCharcoal}{HTML}{2B2826}
\definecolor{BrandRed}{HTML}{B83232}
\definecolor{BrandRedSoft}{HTML}{F9EDED}
\definecolor{BrandMuted}{HTML}{6B6560}
\definecolor{BrandBorder}{HTML}{E8E4DF}
\definecolor{BrandBgSoft}{HTML}{FAF9F7}
\definecolor{InfoBlue}{HTML}{1E5082}
\definecolor{InfoBlueSoft}{HTML}{EDF4FA}
\definecolor{CheckGreen}{HTML}{2D6A4F}
\definecolor{CheckGreenSoft}{HTML}{EDF7F2}

\newcommand{\AuthorName}{Abdellah Raissouni}

% Section styling
\newcommand{\sectionaccent}{\textcolor{BrandRed}{\rule{2.5pt}{1.0em}}}

\titleformat{\section}
  {\fontfamily{phv}\fontseries{b}\selectfont\large\color{BrandCharcoal}}
  {\sectionaccent\hspace{0.55em}\thesection}{0.6em}{#1}
  [\vspace{0.1em}\noindent{\color{BrandBorder}\rule{\textwidth}{0.6pt}}\vspace{0.3em}]
\titlespacing*{\section}{0pt}{1.2em}{0.6em}

\titleformat{\subsection}
  {\fontfamily{phv}\fontseries{b}\selectfont\normalsize\color{BrandCharcoal}}
  {\thesubsection}{0.5em}{#1}
\titlespacing*{\subsection}{0pt}{0.8em}{0.4em}

% Header / footer
\pagestyle{fancy}
\fancyhf{}
\renewcommand{\headrulewidth}{0pt}
\fancyhead[L]{\small\fontfamily{phv}\selectfont\color{BrandMuted}Portail RH Documents}
\fancyhead[R]{\small\fontfamily{phv}\selectfont\color{BrandMuted}Guide de prise en main (v1.0)}
\fancyfoot[C]{\small\fontfamily{phv}\selectfont\color{BrandMuted}\thepage}
\fancypagestyle{plain}{\fancyhf{}\renewcommand{\headrulewidth}{0pt}%
  \fancyfoot[C]{\small\fontfamily{phv}\selectfont\color{BrandMuted}\thepage}}

% Callout boxes
\makeatletter
\newsavebox{\calloutbox}
\newlength{\calloutinnerwidth}
\newcommand{\calloutstart}[2]{%
  \par\medskip\noindent\def\callout@accent{#1}\def\callout@bg{#2}%
  \setlength{\calloutinnerwidth}{\dimexpr\textwidth-28pt\relax}%
  \begin{lrbox}{\calloutbox}\begin{minipage}{\calloutinnerwidth}%
  \RaggedRight
  {\fontfamily{phv}\fontseries{b}\selectfont\color{BrandCharcoal}\callout@title}\par\vspace{0.4em}%
}
\newcommand{\calloutend}{%
  \end{minipage}\end{lrbox}%
  \noindent\begin{tikzpicture}[baseline=(box.base)]%
    \node[fill=\callout@bg, rounded corners=2.5mm, inner xsep=13pt, inner ysep=10pt,
          text width=\calloutinnerwidth, align=left] (box) {\usebox{\calloutbox}};%
    \fill[\callout@accent] (box.south west) rectangle ([xshift=4pt]box.south west |- box.north east);%
  \end{tikzpicture}\par\medskip
}
\newenvironment{infobox}[1]{\def\callout@title{#1}\calloutstart{InfoBlue}{InfoBlueSoft}}{\calloutend}
\newenvironment{warnbox}[1]{\def\callout@title{#1}\calloutstart{BrandRed}{BrandRedSoft}}{\calloutend}
\newenvironment{checkboxbox}[1]{\def\callout@title{#1}\calloutstart{CheckGreen}{CheckGreenSoft}}{\calloutend}
\makeatother

% Helpers
\lstdefinestyle{envblock}{
  basicstyle=\ttfamily\small, breaklines=true, frame=single, framerule=0.4pt,
  rulecolor=\color{BrandBorder}, backgroundcolor=\color{BrandBgSoft},
  xleftmargin=0.5em, columns=fullflexible, keepspaces=true,
  aboveskip=0.6em, belowskip=0.4em,
}
\lstset{style=envblock}

\newcommand{\checkbox}{\tikz[baseline=-0.55ex]{\draw[BrandMuted, line width=0.5pt] (0,0) rectangle (0.52em,0.52em);}}
\newcommand{\mono}[1]{{\ttfamily\nolinkurl{#1}}}
% Bold, brand-colored table header label -- same font family as body text
% (Palatino) so row heights stay consistent between columns; do NOT switch
% \fontfamily per column in a tabularx preamble, it desyncs row baselines
% between columns when combined with \arraystretch.
\newcommand{\tablabel}[1]{\textbf{\textcolor{BrandCharcoal}{#1}}}
\renewcommand{\arraystretch}{1.3}
\setlength{\tabcolsep}{7pt}

% Cover strip
\newcommand{\makecover}{%
  \thispagestyle{empty}
  \begin{tikzpicture}[remember picture, overlay]
    \fill[BrandCharcoal] (current page.north west)
      rectangle ([yshift=-4.2cm]current page.north east);
    \fill[BrandRed]
      ([xshift=-1.6cm,yshift=-4.2cm]current page.north east)
      -- ([xshift=-1.6cm,yshift=-2.4cm]current page.north east)
      -- ([xshift=0cm,  yshift=-4.2cm]current page.north east)
      -- cycle;
    \node[anchor=north west, xshift=2.2cm, yshift=-1.1cm]
      at (current page.north west) {%
      \begin{minipage}{0.75\textwidth}
        {\fontfamily{phv}\fontseries{b}\selectfont\color{white}\fontsize{8}{10}\selectfont
          \MakeUppercase{Amendis / Veolia, Service RH}}\\[0.65em]
        {\fontfamily{phv}\fontseries{b}\selectfont\color{white}\fontsize{22}{26}\selectfont
          Portail RH Documents}\\[0.25em]
        {\fontfamily{phv}\selectfont\color{white!80}\fontsize{12}{16}\selectfont
          Guide de prise en main pour nouvel ingénieur}\\[0.7em]
        {\color{white!65}\fontsize{9.5}{13}\selectfont
          Contexte, fonctionnalités, état d'avancement et démarrage rapide}
      \end{minipage}};
    \node[anchor=south east, xshift=-2.2cm, yshift=1.0cm]
      at (current page.south east) {%
      \begin{minipage}{0.38\textwidth}
        \raggedleft
        {\fontfamily{phv}\selectfont\color{BrandMuted}\small Juillet 2026, version 1.0}\\[0.15em]
        {\fontfamily{phv}\selectfont\color{BrandMuted}\small Contact : \textbf{\AuthorName}}
      \end{minipage}};
  \end{tikzpicture}
  \vspace*{4.6cm}
}

% =============================================================================
\begin{document}
\makecover

\begin{infobox}{Objet de ce document}
  Ce document vous donne une vue d'ensemble du projet avant de plonger dans le
  code : ce qu'il fait, ce qui est déjà construit, la stack technique, comment
  lancer le projet en local, et où trouver le reste de la documentation. Pour
  le détail technique complet (architecture, API, modèle de données,
  conventions), la référence vivante est \mono{doc/TECHNICAL.md}. Lisez-le en
  premier une fois cette vue d'ensemble digérée.
\end{infobox}

{\hypersetup{linkcolor=BrandCharcoal}\tableofcontents}

% -----------------------------------------------------------------------------
\section{Vue d'ensemble}
\label{sec:overview}

Le \textbf{Portail RH Documents} est une application web interne qui donne aux
équipes RH un accès contrôlé aux documents des salariés stockés sur
\textbf{Google Drive}, sans que les agents aient à fouiller manuellement dans
l'arborescence Drive, et sans leur donner un accès plus large que ce dont ils
ont besoin.

\medskip
\textbf{Principe d'architecture central :}

\begin{tabularx}{\textwidth}{@{}p{0.28\textwidth}X@{}}
\toprule
\tablabel{Couche} & \tablabel{Rôle} \\
\midrule
Google Drive & Stockage des fichiers. Reste la source de vérité, les PDF n'y sont jamais copiés ailleurs. \\[0.4em]
PostgreSQL & Index des métadonnées, des utilisateurs et des permissions. C'est sur cette base que la recherche s'exécute, jamais sur Drive en direct. \\[0.4em]
Application Next.js & Authentification, contrôle d'accès, recherche, interface RH/admin. \\
\bottomrule
\end{tabularx}

\medskip
Un document refusé à un utilisateur apparaît dans l'interface \textbf{grisé, sans
lien}. Il n'est jamais masqué en front uniquement : l'accès est toujours vérifié
côté serveur.

% -----------------------------------------------------------------------------
\section{Fonctionnalités principales}

\begin{itemize}[leftmargin=1.4em, itemsep=0.3em, topsep=0.3em]
  \item \textbf{Authentification et comptes} : login, mot de passe oublié, gestion des utilisateurs et des rôles (\mono{ADMIN}, \mono{RH\_MANAGER}, \mono{RH\_AGENT})
  \item \textbf{Recherche documentaire avancée} : recherche plein texte sur toutes les colonnes de la feuille Base (CIN, diplôme, zone, responsable, etc.), filtres par identité, organisation, dates, avec export \textbf{Excel} et export \textbf{ZIP} (arborescence par salarié/catégorie) des résultats
  \item \textbf{Annuaire et dossier salarié} : fiche par matricule avec état civil, documents indexés et liens Drive (si autorisé)
  \item \textbf{Complétude des dossiers} : checklist par salarié des types de documents attendus vs indexés, avec un mécanisme d'attestation manuelle pour les cas non résolubles automatiquement
  \item \textbf{Reporting} : tableau de bord démographie de l'effectif et vue d'ensemble de la complétude
  \item \textbf{Permissions granulaires} : matrice par rôle/utilisateur × catégorie de dossier, plus une portée par salarié (direction/service)
  \item \textbf{Synchronisation Drive} : un Google Apps Script pousse l'arborescence Drive vers le portail. \textbf{Le portail n'appelle jamais l'API Drive lui-même} (voir §\ref{sec:drivesync})
  \item \textbf{Sécurité} : mots de passe chiffrés, limitation de débit, en-têtes de sécurité, webhooks signés HMAC
\end{itemize}

% -----------------------------------------------------------------------------
\section{Stack technique}

\rowcolors{2}{BrandBgSoft}{white}
\begin{tabularx}{\textwidth}{@{}p{0.26\textwidth}X@{}}
\toprule
\rowcolor{white}
\tablabel{Couche} & \tablabel{Choix} \\
\midrule
Framework & Next.js 16 (App Router, Turbopack) \\
Langage & TypeScript \\
UI & Tailwind CSS v4, Shadcn/ui, Framer Motion \\
Formulaires & react-hook-form + zod \\
Auth & NextAuth.js v5 (Credentials, JWT) \\
Base de données & PostgreSQL 16+, Prisma 7 \\
Email & Nodemailer (reset password ; Ethereal en dev) \\
Synchronisation Drive & Google Apps Script (push vers le portail, aucune API Drive côté serveur) \\
Graphiques & recharts (page reporting) \\
\bottomrule
\end{tabularx}

% -----------------------------------------------------------------------------
\section{État d'avancement}

\begin{checkboxbox}{Ce qui est déjà construit}
\begin{tabular}{@{}lp{0.74\textwidth}@{}}
\checkbox & Authentification, gestion des comptes, paramètres utilisateur \\
\checkbox & Synchronisation Drive réelle via Google Apps Script (webhook signé, testée de bout en bout) \\
\checkbox & Recherche documentaire avancée avec exports Excel et ZIP \\
\checkbox & Annuaire, dossier salarié, complétude des dossiers avec attestations manuelles \\
\checkbox & Reporting : démographie de l'effectif et vue de complétude \\
\checkbox & Permissions : matrice par rôle/utilisateur et portée par salarié (direction/service) \\
\checkbox & Sécurité : chiffrement, limitation de débit, en-têtes, webhooks HMAC \\
\end{tabular}
\end{checkboxbox}

% -----------------------------------------------------------------------------
\section{Démarrage rapide}
\label{sec:quickstart}

\begin{enumerate}[leftmargin=1.5em, itemsep=0.3em, topsep=0.2em]
  \item \textbf{Prérequis} : Node.js \textbf{20.9+}, PostgreSQL \textbf{16+} accessible localement, npm.
  \item Cloner le dépôt et installer les dépendances :
\begin{lstlisting}
git clone <url-du-depot>
cd projet-rh-annexe
npm install
\end{lstlisting}
  \item Copier \mono{.env.example} vers \mono{.env} et renseigner au minimum
    \mono{DATABASE\_URL} et \mono{AUTH\_SECRET} (\mono{openssl rand -base64 32}).
    Les variables SMTP et synchronisation Drive peuvent rester vides en local.
  \item Appliquer le schéma Prisma à la base et créer le compte admin :
\begin{lstlisting}
npm run db:setup
\end{lstlisting}
    Sortie réelle observée (schéma déjà à jour, seed idempotent, peut être
    relancé sans risque) :
\begin{lstlisting}[backgroundcolor=\color{BrandCharcoal}, basicstyle=\ttfamily\small\color{white}, rulecolor=\color{BrandCharcoal}]
> projet-rh-annexe@0.1.0 db:setup
> prisma db push && tsx prisma/seed.ts

Loaded Prisma config from prisma.config.ts.
Prisma schema loaded from prisma\schema.prisma.
Datasource "db": PostgreSQL database "rh_annexe", schema "public" at "localhost:5433"

The database is already in sync with the Prisma schema.

Seed OK -- admin@amendis.ma / Admin123!
\end{lstlisting}
  \item Lancer le serveur de développement :
\begin{lstlisting}
npm run dev
\end{lstlisting}
  \item Ouvrir \mono{http://localhost:3000}. Connexion :
    \mono{admin@amendis.ma} / \mono{Admin123!}
\end{enumerate}

\begin{infobox}{Piège fréquent en local}
  Sur certaines machines Windows avec PostgreSQL 18, l'instance locale tourne
  sur le port \textbf{5433} et non le 5432 par défaut. À vérifier avant de
  chercher une erreur de connexion ailleurs (remarquez le port 5433 dans la
  sortie ci-dessus).
\end{infobox}

Pour ce qu'il faut fournir côté IT pour une mise en production (hébergement,
base de données, SMTP, variables d'environnement), voir
\mono{doc/DEPLOYMENT.md}.

% -----------------------------------------------------------------------------
\section{Configurer la synchronisation Drive (Google Apps Script)}
\label{sec:drivesync}

Cette étape n'est \textbf{pas nécessaire pour développer en local} (le point
précédent suffit pour lancer le portail avec une base vide ou déjà peuplée).
Elle concerne la connexion à un vrai Google Drive RH. Référence complète :
\mono{doc/APPS\_SCRIPT\_SYNC\_SETUP.md} ; script source :
\mono{scripts/apps-script/sync.gs} dans le dépôt.

\begin{warnbox}{Aucune API Drive, aucun compte de service}
  Le portail ne détient \textbf{aucune} clé Google, ni compte de service, ni
  client OAuth. Un Google Apps Script (un script qui s'exécute \emph{dans}
  Google, sous l'identité Workspace de l'entreprise) lit le Drive et pousse
  les données vers le portail. Le portail ne fait jamais l'inverse, sauf pour
  l'export ZIP (voir \mono{doc/TECHNICAL.md} §9).
\end{warnbox}

\begin{enumerate}[leftmargin=1.5em, itemsep=0.4em, topsep=0.3em]
  \item \textbf{Localiser le script.} Le code à déployer est dans le dépôt :
    \mono{scripts/apps-script/sync.gs}. Ce fichier n'est \textbf{pas} exécuté
    par Next.js : c'est un fichier source à copier manuellement dans un
    projet Google Apps Script.
  \item \textbf{Créer le projet Apps Script.} Depuis un compte Google
    Workspace de l'entreprise, aller sur \mono{script.google.com} et créer un
    nouveau projet. Plus simple : ouvrir la feuille « Base » (feuille
    Excel/Sheets contenant les salariés) et utiliser
    \mono{Extensions > Apps Script}, ce qui lie automatiquement le script à
    cette feuille.
  \item \textbf{Coller le code.} Copier tout le contenu de
    \mono{scripts/apps-script/sync.gs} dans l'éditeur Apps Script (remplacer
    le fichier \mono{Code.gs} par défaut).
  \item \textbf{Renseigner les propriétés du script.} Dans l'éditeur :
    \mono{Paramètres du projet (icône engrenage) > Propriétés du script >
    Ajouter une propriété du script}. Ajouter les quatre valeurs suivantes :
\begin{lstlisting}
PORTAL_WEBHOOK_URL  = https://<votre-domaine>/api/sync/apps-script
SYNC_SECRET         = <valeur longue aleatoire, ex. openssl rand -base64 32>
ROOT_FOLDER_ID      = <ID du dossier Drive racine, un sous-dossier par matricule>
BASE_SHEET_ID       = <ID de la feuille "Base" (salaries)>
\end{lstlisting}
    L'ID d'un dossier ou d'une feuille Google est la partie de l'URL après
    \mono{/folders/} ou \mono{/d/}.
  \item \textbf{Synchroniser le secret côté portail.} La variable
    d'environnement \mono{DRIVE\_SYNC\_WEBHOOK\_SECRET} du portail doit être
    \textbf{exactement égale} à \mono{SYNC\_SECRET} ci-dessus. Utiliser une
    valeur différente par environnement (dev, staging, prod) : ne jamais
    réutiliser le même secret.
  \item \textbf{Principe du moindre privilège.} Exécuter le script sous un
    compte Workspace qui n'a accès qu'au dossier racine RH et à la feuille
    Base, pas un accès Drive global à tout le domaine.
  \item \textbf{Tester manuellement.} La fonction \mono{testRunSyncManually}
    est fournie dans le script pour lancer une synchronisation à la demande
    depuis l'éditeur (bouton « Exécuter »), avant de programmer le
    déclencheur automatique.
  \item \textbf{Programmer le déclencheur.} Dans l'éditeur Apps Script :
    \mono{Déclencheurs (icône horloge) > Ajouter un déclencheur}, fonction
    \mono{runSync}, type d'événement \mono{Basé sur le temps}, fréquence
    \mono{Toutes les 30 minutes} (ajustable). C'est ce déclencheur, pas un
    bouton dans le portail, qui lance les synchronisations en production.
    L'interface \mono{/admin/sync} du portail est en lecture seule.
\end{enumerate}

\begin{infobox}{Pièges déjà rencontrés en test réel}
  Voir \mono{doc/APPS\_SCRIPT\_SYNC\_SETUP.md} pour le détail, notamment : un
  nom d'onglet \mono{"Base "} avec un espace final fait échouer
  \mono{getSheetByName("Base")} silencieusement ; la ligne 1 de la feuille
  Base doit être vide, les en-têtes en ligne 2 ; et la signature HMAC doit
  être calculée en UTF-8 explicite pour ne pas échouer sur les caractères
  accentués (\mono{Prénom}, \mono{Date d'Entrée}, etc.).
\end{infobox}

% -----------------------------------------------------------------------------
\section{Les branches}
\label{sec:branches}

Le développement se fait sur des branches de fonctionnalité fusionnées dans
\mono{main} une fois stabilisées.

\medskip
\begin{infobox}{main, la branche du projet}
  C'est \textbf{la branche de référence}, celle sur laquelle s'appuyer pour
  démarrer. Elle contient l'architecture actuelle et validée : synchronisation
  Drive via \textbf{Google Apps Script} (le portail ne parle jamais directement
  à l'API Google Drive, voir §\ref{sec:overview} et
  §\ref{sec:drivesync}), recherche PostgreSQL, permissions, exports, et toute
  la documentation à jour dans \mono{doc/}. Sauf indication contraire, c'est
  de là qu'il faut partir.
\end{infobox}

\medskip
Une branche annexe, \mono{feature/local-server-migration}, explore une
direction différente : abandonner Google Drive au profit d'un stockage
\textbf{auto-hébergé} (fichiers sur un serveur local plutôt que sur Drive).
Ce travail n'est \textbf{pas fusionné dans main} et ne représente pas
l'architecture actuelle du projet. À considérer seulement si cette direction
est explicitement remise à l'ordre du jour.

% -----------------------------------------------------------------------------
\section{Documentation}

\rowcolors{2}{BrandBgSoft}{white}
\begin{tabularx}{\textwidth}{@{}p{0.34\textwidth}X@{}}
\toprule
\rowcolor{white}
\tablabel{Fichier} & \tablabel{Contenu} \\
\midrule
\mono{README.md} & Point d'entrée : prérequis, démarrage rapide, scripts npm \\
\mono{doc/TECHNICAL.md} & Référence technique vivante : architecture, API, modèle de données, conventions. À lire en premier. \\
\mono{doc/DEPLOYMENT.md} & Prérequis IT pour la mise en production \\
\mono{doc/APPS\_SCRIPT\_SYNC\_SETUP.md} & Déploiement et sécurité de la synchronisation Google Apps Script \\
\mono{doc/VISUAL\_IDENTITY.md} & Charte graphique et conventions d'interface \\
\mono{doc/FEATURE\_IDEAS.md} & Idées et améliorations en attente \\
\mono{SECURITY.md} & Détails de sécurité et checklist de production \\
\bottomrule
\end{tabularx}

% -----------------------------------------------------------------------------
\section{Pour bien démarrer}

\begin{enumerate}[leftmargin=1.5em, itemsep=0.3em, topsep=0.2em]
  \item Lire \mono{doc/TECHNICAL.md} en entier avant de toucher au code. C'est
    la référence, tenue à jour à chaque changement d'architecture.
  \item Faire tourner le projet en local (§\ref{sec:quickstart}) et se
    connecter avec le compte admin seedé pour explorer l'interface avant de
    lire le code.
  \item Consulter \mono{doc/FEATURE\_IDEAS.md} pour voir ce qui est déjà
    identifié comme prochaine étape.
  \item Travailler à partir de \mono{main} (§\ref{sec:branches}) sauf
    instruction contraire.
\end{enumerate}

\end{document}
